% ===== PRÉAMBULE CANONIQUE — à coller VERBATIM en tête de CHAQUE .tex =====
\documentclass[11pt,a4paper]{article}
\usepackage[utf8]{inputenc}\usepackage[T1]{fontenc}\usepackage{lmodern}
\usepackage[french]{babel}
\usepackage{amsmath,amssymb,mathtools}\usepackage{icomma}
\usepackage[version=4]{mhchem}          % \ce{...} : équations & formules
\usepackage{siunitx}\sisetup{locale=FR} % \qty{}{}, \si{}, \ang{}
\usepackage{chemfig}                    % \chemfig{...} : structures développées
\usepackage{xcolor}\usepackage{colortbl}
\usepackage{tikz}\usepackage{pgfplots}\pgfplotsset{compat=1.18}
\usepackage[most]{tcolorbox}\usepackage{enumitem}\usepackage{array,booktabs}
\usepackage[a4paper,margin=2.2cm]{geometry}
\usetikzlibrary{arrows.meta,calc,angles,quotes,positioning,patterns,backgrounds,shapes.geometric}
\definecolor{primary}{RGB}{222,49,99}\definecolor{primarydark}{RGB}{150,22,55}
\definecolor{defcol}{RGB}{0,114,178}\definecolor{methcol}{RGB}{52,92,148}
\definecolor{excol}{RGB}{0,135,90}\definecolor{attncol}{RGB}{200,40,55}
\tcbset{boxbase/.style={enhanced,breakable,boxrule=0.9pt,arc=2.5pt,
  left=3mm,right=3mm,top=2.3mm,bottom=2.3mm,fonttitle=\bfseries,coltitle=white}}
% --- boîtes TUTORIEL (noms EXACTS, ne pas renommer) ---
\newtcolorbox{cle}[1][]{boxbase,colback=defcol!6,colframe=defcol,
  title={\textsc{À retenir}\ifx\relax#1\relax\relax\else\textmd{ : \itshape #1}\fi}}
\newtcolorbox{methode}[1][]{boxbase,colback=methcol!7,colframe=methcol,sharp corners,
  title={\textsc{Méthode}\ifx\relax#1\relax\relax\else\textmd{ --- \itshape #1}\fi}}
\newtcolorbox{exemplebox}[1][]{boxbase,colback=excol!5,colframe=excol,
  title={\textsc{Exemple}\ifx\relax#1\relax\relax\else\textmd{ : \itshape #1}\fi}}
\newtcolorbox{piege}[1][]{boxbase,colback=attncol!6,colframe=attncol,
  title={\textsc{Piège}\ifx\relax#1\relax\relax\else\textmd{ : \itshape #1}\fi}}
\newtcolorbox{remarque}[1][]{boxbase,colback=black!4,colframe=black!45,
  title={\textsc{Remarque}\ifx\relax#1\relax\relax\else\textmd{ : \itshape #1}\fi}}
% --- boîtes EXERCICE / CORRIGÉ (noms EXACTS, auto-comptées) ---
\newtcolorbox[auto counter]{exo}[1][]{boxbase,colback=primary!4,colframe=primary,
  title={\textsc{Exercice}~\thetcbcounter\ifx\relax#1\relax\relax\else\textmd{ --- \itshape #1}\fi}}
\newtcolorbox[auto counter]{corrige}[1][]{boxbase,colback=excol!4,colframe=excol,
  title={\textsc{Corrigé}~\thetcbcounter\ifx\relax#1\relax\relax\else\textmd{ --- \itshape #1}\fi}}
\newcommand{\R}{\mathbb{R}}\newcommand{\N}{\mathbb{N}}\newcommand{\Z}{\mathbb{Z}}\newcommand{\ds}{\displaystyle}
% (un \usepackage{fancyhdr}+en-tête est possible ; il est ignoré côté web)
% =========================================================================
\begin{document}

\begin{exo}[deux enceintes à la même température]
Une enceinte de $V_1=\qty{5.0}{\litre}$ contient $\qty{16}{\gram}$ de méthane \ce{CH4}
($M=\qty{16}{\gram\per\mole}$) sous $p_1=\qty{0.4}{atm}$. Une seconde enceinte de $V_2=\qty{10.0}{\litre}$
contient $\qty{16}{\gram}$ de dioxygène \ce{O2} ($M=\qty{32}{\gram\per\mole}$).
\begin{enumerate}[label=\alph*)]
  \item Calcule la quantité de matière $n_1$ de \ce{CH4}.
  \item Calcule la quantité de matière $n_2$ de \ce{O2}.
  \item En écrivant que la température est \emph{la même} dans les deux enceintes, exprime la pression
        $p_2$ en fonction de $p_1$, des volumes et des quantités de matière.
  \item Calcule $p_2$.
\end{enumerate}
\end{exo}

\begin{exo}[combustion du méthane : volumes de gaz aux CNTP]
On brûle complètement $\qty{8.0}{\gram}$ de méthane selon
\[ \ce{CH4 + 2 O2 -> CO2 + 2 H2O}. \]
On donne $M(\ce{CH4})=\qty{16}{\gram\per\mole}$ ; $V_{\mathrm{m}}=\qty{22.4}{\litre\per\mole}$ aux CNTP.
\begin{enumerate}[label=\alph*)]
  \item Calcule la quantité de matière de \ce{CH4} brûlé.
  \item Déduis la quantité de \ce{CO2} formé.
  \item Quel volume de \ce{CO2} cela représente-t-il aux CNTP ?
  \item Quel volume de \ce{O2} (aux CNTP) a été consommé ?
\end{enumerate}
\end{exo}

\begin{exo}[gaz d'un airbag : de la masse au volume]
Le gonflage d'un airbag repose sur la décomposition de l'azoture de sodium :
\[ \ce{2 NaN3(s) -> 2 Na(s) + 3 N2(g)}. \]
On décompose $\qty{6.5}{\gram}$ de \ce{NaN3} ($M=\qty{65}{\gram\per\mole}$) ;
$V_{\mathrm{m}}=\qty{22.4}{\litre\per\mole}$ aux CNTP.
\begin{enumerate}[label=\alph*)]
  \item Calcule la quantité de matière de \ce{NaN3}.
  \item Déduis la quantité de \ce{N2} produite.
  \item Quel volume de \ce{N2} obtient-on aux CNTP ?
\end{enumerate}
\end{exo}

\begin{exo}[transformation d'un gaz depuis les CNTP]
Un système contient $\qty{0.20}{\gram}$ de dihydrogène \ce{H2} ($M=\qty{2}{\gram\per\mole}$) dans les
conditions normales. On l'amène ensuite à $T=\qty{300}{\kelvin}$ sous $p=\qty{1.5}{atm}$ (on prend
$R=\num{0,082}$).
\begin{enumerate}[label=\alph*)]
  \item Calcule la quantité de matière $n$ de \ce{H2}.
  \item Calcule le volume $V_1$ occupé aux CNTP.
  \item Calcule le volume $V_2$ dans l'état final ($\qty{300}{\kelvin}$, $\qty{1.5}{atm}$).
  \item En déduis la variation de volume $\Delta V=V_2-V_1$.
\end{enumerate}
\end{exo}

\end{document}
